\begin{pastexam}[label={px:ust-114-ee-ds-24}]{臺聯大電機類 114 學年度 第 24 題}
（複選；5 分，答錯倒扣 1 分，倒扣至本大題 0 分為止）

We have an open addressing hash table. Please simulate the hash table's behavior storing integers
given the following conditions:\\
$\checkmark$The hash table size is 9.\\
$\checkmark$It uses quadratic probing for collision resolution.\\
$\checkmark$The hash function uses the int value (plus any probing needed) mod the size of the table as
described in the following.\\
Hashing+Probing: ((k mod m) + i$^2$) mod m, k is key, m is the size of the table, probing from i = 0,
1, 2,\ldots\\
What values will be in the hash table after the following sequence of insertions?\\
18, 16, 10, 7, 26

\begin{center}
\begin{tabular}{|*{9}{p{1.1cm}|}}
\hline
[0] & [1] & [2] & [3] & [4] & [5] & [6] & [7] & [8] \\ \hline
 & & & & & & & & \\ \hline
\end{tabular}
\end{center}

Which of the following statements are correct. (5\%) multiple anwsers
\begin{enumerate}[label=(\Alph*)]
\item 18 is stored at [0]
\item 10 is stored at [1]
\item 26 is stored at [8]
\item 7 is stored at [7]
\item 16 is stored at [7]
\end{enumerate}
\end{pastexam}

\begin{pastsoln}
答案是 (A)(B)(E)。

$m = 9$，$h(k, i) = ((k \bmod 9) + i^2) \bmod 9$（資結書第~12.3~節的二次探測）。 逐一插入：
\begin{center}
\small
\begin{tabular}{cccl}
\toprule
$k$ & $k \bmod 9$ & 依序試的格子 & 結果 \\
\midrule
18 & 0 & 0 & 放 [0] \\
16 & 7 & 7 & 放 [7] \\
10 & 1 & 1 & 放 [1] \\
7 & 7 & 7（有 16）、$7 + 1 = 8$ & 放 [8] \\
26 & 8 & 8（有 7）、$(8 + 1) \bmod 9 = 0$（有 18）、$(8 + 4) \bmod 9 = 3$ & 放 [3] \\
\bottomrule
\end{tabular}
\end{center}
最後的表：
\begin{center}
\begin{tabular}{|*{9}{c|}}
\hline
[0] & [1] & [2] & [3] & [4] & [5] & [6] & [7] & [8] \\ \hline
18 & 10 & & 26 & & & & 16 & 7 \\ \hline
\end{tabular}
\end{center}
\begin{itemize}
\item (A) 18 在 [0]：對。
\item (B) 10 在 [1]：對。
\item (C) 26 在 [8]：[8] 已被 7 佔用，26 最後在 [3]。 錯。
\item (D) 7 在 [7]：[7] 先被 16 佔用，7 在 [8]。 錯。
\item (E) 16 在 [7]：對。
\end{itemize}
\end{pastsoln}
